\subsection{证明细节：选择对应、最大值与比较静态}

\begin{proposition}[选择对应的集合值性质]
设 $C(\theta)=\argmax_{x\in B(\theta)}u(x)$。若 $B(\theta)$ 非空紧、$u$ 连续，则 $C(\theta)$ 非空闭，并且 $C(\theta)\subseteq B(\theta)$。
\end{proposition}
\begin{proof}
包含关系是 argmax 定义的直接结果。非空性由最大值定理得到。若 $x^n\in C(\theta)$ 且 $x^n\to x$，因为 $B(\theta)$ 闭，$x\in B(\theta)$；又因 $u(x^n)=M$ 且 $u$ 连续，$u(x)=M$，故 $x\in C(\theta)$。因此选择对应是闭集，而不是任意的点集。
\end{proof}

\begin{example}[线性效用的整段最优]
令 $u(x)=x_1+2x_2$，预算线为 $p_1x_1+p_2x_2=m$。若 $p_1/p_2=1/2$，预算线与效用等值线平行，所有预算边界点都是最优点；若 $p_1/p_2<1/2$，单位支出带来的商品一效用更高，最优为商品一角点；若 $p_1/p_2>1/2$，最优为商品二角点。三种情形分别对应整段、左端点和右端点，不能仅用一个“切点”公式概括。
\end{example}

\begin{proposition}[严格拟凹性与唯一性]
若 $B$ 为凸集、$u$ 严格拟凹，则 $\argmax_Bu$ 至多包含一个点。
\end{proposition}
\begin{proof}
反设 $x\ne y$ 都是最大点，记最大值为 $M$。凸性给出 $z=\lambda x+(1-\lambda)y\in B$，其中 $0<\lambda<1$。严格拟凹性给出
$u(z)>\min\{u(x),u(y)\}=M$，与 $M$ 是最大值矛盾。因此最优点至多一个。
\end{proof}

\subsection{参数变化的逐步判别}
预算集 $B(p,m)$ 随收入 $m$ 增加而扩张，随价格上升而收缩，但最优集合的变化还取决于效用函数。若 $u$ 严格单调，最优点位于边界；若 $u$ 存在饱和点，收入增加可能不改变选择。比较静态必须同时写出约束变化和目标函数结构，不能从集合扩张直接推出每个消费量上升。

\begin{example}[收入扩张与非正规品]
若效用为 $u(x_1,x_2)=x_1-(x_2-2)^2$，收入增加首先使 $x_2$ 向饱和点 $2$ 调整，超过该点后继续增加收入并不必然增加 $x_2$。这说明“收入上升、需求上升”是正规品假设下的结论，而非预算集扩张的逻辑必然。
\end{example}

\begin{proposition}[包络公式的直接证明]
令 $v(\theta)=u(x(\theta),\theta)$，且 $x(\theta)$ 为内点唯一最优解、满足一阶条件 $u_x=0$。则
$v_{\theta}=u_{\theta}(x(\theta),\theta)$。
\end{proposition}
\begin{proof}
链式法则给出 $v_{\theta}=u_x(x(\theta),\theta)x_{\theta}+u_{\theta}(x(\theta),\theta)$。最优一阶条件令第一项为零，故只剩直接参数效应。约束问题中把约束乘子写入拉格朗日函数后，同样的消去机制给出价格和收入的包络导数。
\end{proof}

\subsection{模型边界的三个反例}
\begin{enumerate}
\item 若 $p_i=0$，预算集沿商品 $i$ 方向可能无界，紧性论证失效。
\item 若 $u$ 不连续，闭预算集也不能保证最大值达到；跳跃效用可能只有上确界。
\item 若可行集不凸，一阶条件即使存在，也可能只识别某一局部峰值。
\end{enumerate}
上述反例不是新增模型，而是对本章既有存在性、唯一性和一阶条件结论的边界核查。写作时应把每个定理的假设与结论逐项对应，避免把一个结论的假设带入另一个结论。

\subsection{同一选择问题的三种边界情形}
设预算集为 $B(p,m)=\{x\in\mathbb R_+^2:p_1x_1+p_2x_2\le m\}$。对线性效用
$u(x)=a_1x_1+a_2x_2$，比较单位支出的效用收益：
\[
\rho_1=\frac{a_1}{p_1},\qquad \rho_2=\frac{a_2}{p_2}.
\]
若 $\rho_1>\rho_2$，最优集合退化为商品一角点；若 $\rho_1<\rho_2$，最优集合退化为商品二角点。只有 $\rho_1=\rho_2$ 时，预算边界上的每一点都给出相同效用，因而
\[
C(p,m)=\{x\ge0:p\cdot x=m\}.
\]
这三个区域说明，角点、整段和唯一内点并不是互相矛盾的结论，而是同一最大化问题在不同参数区域的完整划分。

\begin{proposition}[财富耗尽的条件]
若 $u$ 对每个商品都严格单调、$p\gg0$ 且 $m>0$，则任意 $x^*\in C(p,m)$ 都满足 $p\cdot x^*=m$。
\end{proposition}
\begin{proof}
反设 $p\cdot x^*<m$。由于 $p_i>0$，可取充分小的 $\delta>0$，使得
$p\cdot(x^*+\delta e_i)\le m$。严格单调性给出
$u(x^*+\delta e_i)>u(x^*)$，这与 $x^*$ 已达到可行集上的最大值矛盾。因此最优点必须位于预算边界。若效用只满足弱单调，财富耗尽则不再自动成立。
\end{proof}

\subsection{从候选点到全局最优：一项完整核对}
设 $B=\{x\ge0:p\cdot x\le m\}$，且 $u$ 连续、严格递增并严格拟凹。首先由 $p\gg0$ 验证 $B$ 有界，由闭性验证 $B$ 闭，故 $B$ 紧；最大值定理给出最优解存在。其次，严格递增性排除预算内部点，故 $p\cdot x^*=m$。最后，严格拟凹性排除两个不同的最大点，因而最优解唯一。

这一链条中每一步承担不同任务：紧性只负责“能取到最大值”，严格单调只负责“财富耗尽”，严格拟凹只负责“不能有两个最优点”。若删除其中任一条件，结论都必须相应减弱。比如连续但不严格拟凹的线性效用可以产生整段最优；严格拟凹但预算集不紧时，可能只有上确界而没有最大值。

\begin{example}[两项假设的逐项删去]
取 $u(x)=x_1+x_2$，它连续但不是严格拟凹，预算线上所有点都最优；若预算集改为开集 $\{x\ge0:p\cdot x<m\}$，即使 $u$ 连续且严格递增，也只能不断逼近边界，最大值可能不存在。这两个例子分别对应唯一性与存在性的失败。
\end{example}

\subsection{候选点、边界点与全局比较的统一算法}
对于本章反复出现的约束最大化问题，可以把证明拆成四个互不替代的步骤。第一步写出可行集并逐项检查非空、闭性和有界性；第二步利用连续性得到最大值存在；第三步根据单调性判断最优解是否耗尽约束；第四步才使用一阶条件、Kuhn--Tucker 条件或无差异曲线进行定位。最后必须将候选点与全部边界点比较，才能把局部条件升级为全局结论。

例如，对 $u(x)=\alpha\log x_1+(1-\alpha)\log x_2$，定义域要求 $x_1,x_2>0$。在正象限内，拉格朗日一阶条件给出
\[
\frac{\alpha}{x_1}=\lambda p_1,\qquad
\frac{1-\alpha}{x_2}=\lambda p_2.
\]
两式相除得到 $x_1/x_2=\alpha p_2/[(1-\alpha)p_1]$，再代入预算等式得到候选束。由于对数效用严格凹，任意满足一阶条件的候选束都是唯一全局最优；这里“严格凹”承担的是充分性和唯一性，而不是存在性本身。

若把效用换成线性函数，上述一阶条件可能没有内点解，必须比较两个角点和预算线整段。若把预算集换成不闭的开集，连续性也不再足以保证最大值达到。由此可见，证明顺序不是形式上的排版要求，而是防止把候选解误写成最优解的逻辑防线。

\subsection{一阶条件的充分性为何需要曲率条件}
设 $x^*$ 是内点候选解，满足 $\nabla u(x^*)=\lambda p$。这一等式只说明沿预算超平面的局部可行方向，效用的一阶变化为零，并不能单独排除另一个更高的可行点。若 $u$ 凹，则任意 $y\in B(p,m)$ 满足
\[
 u(y)\le u(x^*)+\nabla u(x^*)\cdot(y-x^*)
 =u(x^*)+\lambda p\cdot(y-x^*)\le u(x^*),
\]
因此候选点是全局最优。若 $u$ 严格凹，最后一个不等式在 $y\ne x^*$ 时严格，故最优唯一。这个两行不等式正是“局部一阶条件加全局曲率”升级为全局结论的关键。

若只有拟凹性而没有可微性，则不能直接写梯度条件，但仍可用上水平集的凸性和支持超平面描述最优点；若有可微性而没有凹性，则一阶条件只给出驻点，必须回到边界比较或二阶条件。因而“写出一阶条件”与“证明最优性”在逻辑上始终是两件事。

\subsection{本章证明链的回顾与应用边界}
从方案空间到选择对应，完整论证依次经过集合的定义、可行性的限制、最大值的存在、最优集合的识别以及参数变化的比较。预算集的紧性保证“最大值能够取到”，连续性把集合性质传递给目标函数，单调性决定边界位置，严格拟凹性决定唯一性，而微分条件只是在这些前提成立后提供计算工具。任何一步被省略，结论都只能保持在较弱层次。

因此，在阅读或书写消费者问题时，应同时记录“结论—假设—证明工具”三列。例如，存在性对应紧性与连续性；财富耗尽对应严格单调性与正价格；唯一性对应严格拟凹性；需求导数对应内点、可微性和隐函数定理。这个核对表既能解释定理为何成立，也能在参数改变时迅速指出哪一条结论需要重新证明。

\begin{example}[同一模型中的假设替换]
若把严格单调替换为弱单调，预算耗尽结论必须改写为“至少存在一个边界最优解”，不能再断言每个最优解都在边界；若把严格拟凹替换为普通拟凹，存在性仍可保留，但唯一性消失；若把闭预算集替换为开预算集，连续性仍可保留，但最大值定理不再适用。逐项替换说明，假设不是装饰性的前置条件，而是分别支撑不同结论的逻辑支柱。
\end{example}

\subsection{参数、集合与选择的三层区分}
设参数从 $\theta$ 变为 $\theta'$，首先要比较的是可行集 $B(\theta)$ 与 $B(\theta')$；其次才比较目标函数在两个集合上的最大值；最后比较选择对应 $C(\theta)$ 与 $C(\theta')$。即使 $B(\theta)\subseteq B(\theta')$，也只能说明可行机会扩大，不能单独推出每一坐标需求增加，更不能推出福利必然增加，除非再加入单调性、补偿或规范评价的额外条件。

在唯一最优的情形，选择对应可写成函数 $x(\theta)$，可以讨论连续性和导数；在多重最优的情形，必须保留集合值记号 $C(\theta)$，并使用上半连续、下半连续或选择函数来描述变化。把集合值对应强行写成单值函数，会掩盖线性效用、完全替代和约束退化时的真实结构。该区分也是本章从集合语言出发、最后回到模型价值评价的闭环。
